% Part: many-valued-logic
% Chapter: three-valued-logics
% Section: kleene

\documentclass[../../../include/open-logic-section]{subfiles}

\begin{document}

\olfileid[id]{mvl}{thr}{skl}

\olsection{Logika Kleene}

Stephen Kleene memperkenalkan dua logika bernilai tiga yang dimotivasi oleh
suatu pandangan bahwa nilai kebenaran merupakan hasil prosedur komputasi:
suatu prosedur mungkin menghasilkan $\True$ atau~$\False$, tetapi mungkin pula
tidak pernah berhenti. Dalam kasus terakhir, nilai kebenaran yang bersesuaian
tidak terdefinisi dan direpresentasikan oleh nilai kebenaran~$\Undef$.

Untuk menghitung negasi suatu proposisi~$!A$, mula-mula kita menghitung nilai
dari~$!A$, kemudian mengembalikan kebalikan hasil tersebut. Jika komputasi
$!A$ tidak berhenti, keseluruhan prosedur itu pun tidak berhenti; jadi negasi
$\Undef$ adalah~$\Undef$.

Untuk menghitung suatu konjungsi $!A\land !B$, terdapat dua pilihan. Kita dapat
menghitung~$!A$ terlebih dahulu, lalu~$!B$; hasilnya adalah~$\True$ jika kedua
komputasi menghasilkan~$\True$, dan~$\False$ dalam kasus lainnya. Jika salah
satu komputasi tidak berhenti, keseluruhan prosedur pun tidak berhenti. Jadi,
dalam cara ini, jika salah satu konjung tidak terdefinisi, konjungsinya juga
tidak terdefinisi. Hal yang sama berlaku bagi disjungsi.

Akan tetapi, jika kita dapat mengevaluasi $!A$ dan~$!B$ secara paralel, kita
dapat memperoleh hasil yang lebih baik. Jika salah satu dari kedua prosedur
berhenti dan mengembalikan~$\False$, kita dapat berhenti karena jawabannya
pasti salah. Dengan demikian, konjungsi dengan satu konjung salah bernilai
salah, meskipun konjung lainnya tidak terdefinisi. Serupa dengan itu, ketika
menghitung suatu disjungsi secara paralel, kita dapat berhenti segera setelah
prosedur bagi salah satu disjung mengembalikan benar: disjungsi tersebut pasti
benar. Jadi, dalam cara ini kita dapat mengetahui hasil suatu klaim majemuk
meskipun salah satu komponennya tidak terdefinisi. Dalam penafsiran ini,
$\Undef$ dapat dibaca sebagai ``tidak diketahui'', bukan ``tidak
terdefinisi''.

Kedua penafsiran tersebut menghasilkan logika Kleene kuat dan lemah.
Implikasi didefinisikan sebagai padanan dari~$\lnot !A\lor !B$.

\begin{defn}
\emph{Logika Kleene kuat}~$\LogKs$ didefinisikan dengan menggunakan matriks
berikut:
\begin{enumerate}
  \item Bahasa proposisional baku~$\Lang L_0$ dengan $\lnot$, $\land$,
  $\lor$, $\lif$.
  \item Himpunan nilai kebenaran~$V=\{\True,\Undef,\False\}$.
  \item $\True$ merupakan satu-satunya nilai yang ditetapkan, yakni
  $V^+=\{\True\}$.
  \item Fungsi kebenaran diberikan oleh tabel-tabel berikut:
  \begin{center}
    \begin{tabular}{c|c}
      $\tf{\lnot}$ & \\
      \hline
      $\True$ & $\False$ \\
      $\Undef$ & $\Undef$ \\
      $\False$ & $\True$
    \end{tabular}
    \quad
    \begin{tabular}{c|ccc}
      $\tf{\land}[\LogKs]$ & $\True$ & $\Undef$ & $\False$ \\
      \hline
      $\True$ & $\True$ & $\Undef$ & $\False$ \\
      $\Undef$ & $\Undef$ & $\Undef$ & $\False$\\
      $\False$ & $\False$ & $\False$ & $\False$
    \end{tabular}
    \\[2ex]
    \begin{tabular}{c|ccc}
      $\tf{\lor}[\LogKs]$ & $\True$ & $\Undef$ & $\False$ \\
      \hline
      $\True$ & $\True$ & $\True$ & $\True$ \\
      $\Undef$ & $\True$ & $\Undef$ & $\Undef$ \\
      $\False$ & $\True$ & $\Undef$ & $\False$
    \end{tabular}
    \quad
    \begin{tabular}{c|ccc}
      $\tf{\lif}[\LogKs]$ & $\True$ & $\Undef$ & $\False$ \\
      \hline
      $\True$ & $\True$ & $\Undef$ & $\False$ \\
      $\Undef$ & $\True$ & $\Undef$ & $\Undef$  \\
      $\False$ & $\True$ & $\True$ & $\True$
    \end{tabular}
  \end{center}
\end{enumerate}
\end{defn}

\begin{defn}
\emph{Logika Kleene lemah}~$\LogKw$ didefinisikan dengan menggunakan matriks
berikut:
\begin{enumerate}
  \item Bahasa proposisional baku~$\Lang L_0$ dengan $\lnot$, $\land$,
  $\lor$, $\lif$.
  \item Himpunan nilai kebenaran~$V=\{\True,\Undef,\False\}$.
  \item $\True$ merupakan satu-satunya nilai yang ditetapkan, yakni
  $V^+=\{\True\}$.
  \item Fungsi kebenaran diberikan oleh tabel-tabel berikut:
  \begin{center}
    \begin{tabular}{c|c}
      $\tf{\lnot}$ & \\
      \hline
      $\True$ & $\False$ \\
      $\Undef$ & $\Undef$ \\
      $\False$ & $\True$
    \end{tabular}
    \quad
    \begin{tabular}{c|ccc}
      $\tf{\land}[\LogKw]$ & $\True$ & $\Undef$ & $\False$ \\
      \hline
      $\True$ & $\True$ & $\Undef$ & $\False$ \\
      $\Undef$ & $\Undef$ & $\Undef$ & $\Undef$\\
      $\False$ & $\False$ & $\Undef$ & $\False$
    \end{tabular}
    \\[2ex]
    \begin{tabular}{c|ccc}
      $\tf{\lor}[\LogKw]$ & $\True$ & $\Undef$ & $\False$ \\
      \hline
      $\True$ & $\True$ & $\Undef$ & $\True$ \\
      $\Undef$ & $\Undef$ & $\Undef$ & $\Undef$ \\
      $\False$ & $\True$ & $\Undef$ & $\False$
    \end{tabular}
    \quad
    \begin{tabular}{c|ccc}
      $\tf{\lif}[\LogKw]$ & $\True$ & $\Undef$ & $\False$ \\
      \hline
      $\True$ & $\True$ & $\Undef$ & $\False$ \\
      $\Undef$ & $\Undef$ & $\Undef$ & $\Undef$  \\
      $\False$ & $\True$ & $\Undef$ & $\True$
    \end{tabular}
  \end{center}
\end{enumerate}
\end{defn}

\begin{prop}
  $\LogKs$ dan~$\LogKw$ tidak mempunyai tautologi.
\end{prop}

\begin{proof}
  Jika $\pAssign v(p)=\Undef$ bagi semua !!{propositional variable}s~$p$,
  maka setiap formula~$!A$ mempunyai nilai kebenaran
  $\pValue v(!A)=\Undef$, sebab
  \[
    \tf{\lnot}(\Undef)=\tf{\lor}(\Undef,\Undef)=\tf{\land}(\Undef,
  \Undef)=\tf{\lif}(\Undef,\Undef)=\Undef
  \]
  dalam kedua logika. Karena $\Undef\notin V^+$ baik bagi $\LogKs$ maupun
  bagi~$\LogKw$, pada !!{valuation} ini nilai~$!A$ tidak ditetapkan.
\end{proof}

Meskipun logika Kleene lemah dan kuat sama-sama tidak mempunyai tautologi,
keduanya mempunyai relasi konsekuensi yang tidak trivial.

\begin{prob}
  Manakah di antara relasi-relasi berikut yang berlaku dalam (a)~logika
  Kleene kuat dan (b)~logika Kleene lemah? Berikan tabel kebenaran bagi
  masing-masing relasi.
  \begin{enumerate}
    \item $p,p\lif q\Entails q$
    \item $p\lor q,\lnot p\Entails q$
    \item $p\land q\Entails p$
    \item $p\Entails p\land p$
    \item $p\Entails p\lor q$
  \end{enumerate}
\end{prob}

Dmitry Bochvar menafsirkan $\Undef$ sebagai ``tidak bermakna'' dan berusaha
menggunakannya untuk menyelesaikan paradoks seperti paradoks Pembohong dengan
menetapkan bahwa kalimat paradoksal bernilai~$\Undef$. Ia memperkenalkan suatu
logika yang pada dasarnya merupakan logika Kleene lemah yang diperluas dengan
konektif tambahan, dua di antaranya adalah ``negasi eksternal'' dan operator
``tidak terdefinisi'':
\begin{center}
  \begin{tabular}{c|c}
    $\tf{\sim}$ & \\
    \hline
    $\True$ & $\False$ \\
    $\Undef$ & $\True$ \\
    $\False$ & $\True$
  \end{tabular}
\quad
  \begin{tabular}{c|c}
    $\tf{+}$ & \\
    \hline
    $\True$ & $\False$ \\
    $\Undef$ & $\True$ \\
    $\False$ & $\False$
  \end{tabular}
\end{center}

\begin{prob}
  Dapatkah Anda mendefinisikan $\sim$ dalam logika Bochvar dengan menggunakan
  $\lnot$ dan~$+$; yakni, dapatkah Anda menemukan !!a{formula} yang hanya
  memuat !!{propositional variable}~$p$, tidak memuat~$\sim$, dan selalu
  mempunyai nilai kebenaran yang sama dengan~$\mathord{\sim}p$? Berikan tabel
  kebenaran untuk menunjukkan bahwa jawaban Anda benar.
\end{prob}

\end{document}
